\documentclass[11pt]{amsart}
\usepackage{amsmath,amssymb,amsthm}
\usepackage[hidelinks]{hyperref}

\newtheorem{theorem}{Theorem}
\newtheorem{lemma}{Lemma}

\newcommand{\N}{\mathbb N}
\newcommand{\F}{\mathcal F}
\newcommand{\A}{\mathcal A}
\newcommand{\B}{\mathcal B}
\newcommand{\Lp}{L^p}
\newcommand{\psip}{\psi'}

\title{The two-subalgebra sharpness case for a marginal-subspace conjecture}
\author{agent\_lane\_04}
\date{2026-06-29}

\begin{document}
\maketitle

\section*{Source and target}

The source is Ivan Feshchenko, \emph{When is the sum of complemented
subspaces complemented?}, arXiv:1606.08048.  In Section 3.4, after proving
that \(r(E)<1\) is sufficient for the centered marginal subspaces
\[
  L^p_0(\F_1)+\cdots+L^p_0(\F_n)
\]
to be complemented, the paper asks whether this condition is sharp and states
a conjecture: if \(E=(e_{ij})\) is symmetric, has zero diagonal, nonnegative
off-diagonal entries, and \(r(E)=1\), then there should be a probability
space and sub-\(\sigma\)-algebras \(\F_1,\ldots,\F_n\) with
\[
  \psip(\F_i,\F_j)=1-e_{ij}
\]
such that the centered marginal sum is not closed in \(L^p\) for arbitrary
\(p\in[1,\infty]\).

This packet proves the conjecture for \(n=2\).  For \(n=2\), the only such
matrix is
\[
  E=\begin{pmatrix}0&1\\1&0\end{pmatrix}.
\]

\section*{Definitions}

For sub-\(\sigma\)-algebras \(\A,\B\) of a probability space
\((\Omega,\F,\mu)\), Feshchenko uses
\[
  \psip(\A,\B)
  =
  \inf\left\{
    \frac{\mu(A\cap B)}{\mu(A)\mu(B)}:
    A\in\A,\ B\in\B,\ \mu(A)>0,\ \mu(B)>0
  \right\}.
\]
Also \(L^p_0(\A)\) denotes the \(\A\)-measurable \(L^p\)-functions of mean
zero.

\section*{Main result}

\begin{theorem}
There is a probability space \((\Omega,\F,\mu)\) and two sub-\(\sigma\)-algebras
\(\A,\B\subset\F\) such that
\[
  \psip(\A,\B)=0
\]
and, for every \(p\in[1,\infty]\),
\[
  L^p_0(\A)+L^p_0(\B)
\]
is not closed in \(L^p(\Omega,\F,\mu)\).
Consequently Feshchenko's Section 3.4 conjecture is true for \(n=2\).
\end{theorem}

\section*{Proof intuition}

The sample space is the edge set of an infinite bipartite path.  Row-measurable
functions are potentials on one side of the path, and column-measurable
functions are potentials on the other side.  A row potential \(a_i\) and a
column potential \(-a_j\) cancel exactly on the diagonal edges \((i,j)=(m,m)\)
and cancel up to the increment \(a_{m+1}-a_m\) on the adjacent edges
\((m+1,m)\).  Long, slowly varying triangular potentials therefore give
row/column functions whose sum is very small although neither centered
function is small.  This rules out closedness of the sum.

\section*{Construction}

Choose odd integers \(L_k\ge 5\) with \(L_k\to\infty\), and partition
\(\N\) into consecutive intervals
\[
  J_k=\{s_k,s_k+1,\ldots,s_k+L_k-1\}.
\]
Put
\[
  q_m=\frac{2^{-k}}{L_k}\qquad(m\in J_k).
\]
Then \(\sum_{m\ge1}q_m=1\).

Let
\[
  \Omega=\{a_m,b_m:m\in\N\},
\]
where \(a_m\) is the edge joining row \(m\) to column \(m\), and \(b_m\) is the
edge joining row \(m+1\) to column \(m\).  Define
\[
  \mu(\{a_m\})=\mu(\{b_m\})=\frac{q_m}{2}.
\]
Let \(\A\) be the \(\sigma\)-algebra generated by the row map
\[
  \rho(a_m)=m,\qquad \rho(b_m)=m+1,
\]
and let \(\B\) be the \(\sigma\)-algebra generated by the column map
\[
  \kappa(a_m)=\kappa(b_m)=m.
\]

\begin{lemma}
\(\psip(\A,\B)=0\), and \(L^p_0(\A)\cap L^p_0(\B)=\{0\}\) for every
\(p\in[1,\infty]\).
\end{lemma}

\begin{proof}
For instance, the row atom \(\rho^{-1}(\{1\})\) and the column atom
\(\kappa^{-1}(\{3\})\) have positive measure and empty intersection.  Hence
\(\psip(\A,\B)=0\).

If an \(\A\)-measurable function \(f\) and a \(\B\)-measurable function \(g\)
agree almost everywhere, write their row and column values as \(\alpha_i\)
and \(\beta_j\).  Since every edge has positive measure, the edge \(a_m\)
gives \(\alpha_m=\beta_m\), while the edge \(b_m\) gives
\(\alpha_{m+1}=\beta_m\).  Thus
\[
  \alpha_m=\beta_m=\alpha_{m+1}
\]
for every \(m\), so all row and column values are equal to one constant.
If the common function has mean zero, that constant is zero.
\end{proof}

\section*{Failure of closedness}

For \(m\in J_k\), let \(c_k\) be the midpoint of \(J_k\) and put
\[
  \theta_k(m)=1-\frac{|m-c_k|}{(L_k-1)/2};
\]
put \(\theta_k(m)=0\) for \(m\notin J_k\).  Then
\[
  0\le \theta_k\le1,\qquad
  \max_m\theta_k(m)=1,\qquad
  |\theta_k(m+1)-\theta_k(m)|\le \frac{2}{L_k-1}.
\]
Define bounded functions
\[
  u_k(\omega)=\theta_k(\rho(\omega)),\qquad
  v_k(\omega)=-\theta_k(\kappa(\omega)).
\]
Thus \(u_k\) is \(\A\)-measurable and \(v_k\) is \(\B\)-measurable.  On
\(a_m\) one has \(u_k+v_k=0\), and on \(b_m\) one has
\[
  (u_k+v_k)(b_m)=\theta_k(m+1)-\theta_k(m).
\]
Hence
\[
  \|u_k+v_k\|_\infty\le \frac{2}{L_k-1}. \tag{1}
\]

Let
\[
  U_k=u_k-\int u_k\,d\mu,\qquad
  V_k=v_k-\int v_k\,d\mu.
\]
Then \(U_k\in L^p_0(\A)\) and \(V_k\in L^p_0(\B)\) for every \(p\).  Since
\[
  U_k+V_k=(u_k+v_k)-\int (u_k+v_k)\,d\mu,
\]
(1) gives
\[
  \|U_k+V_k\|_\infty\le \frac{4}{L_k-1}. \tag{2}
\]

For finite \(p\), let \(W_k=2^{-k}\), the total mass of the columns in
\(J_k\).  A fixed positive fraction of the column atoms in \(J_k\) have
\(\theta_k\ge 1/2\), while \(\int |v_k|\,d\mu\le W_k\to0\).  Therefore, for
each \(p<\infty\), there is a constant \(c_p>0\), independent of \(k\), such
that for all large \(k\)
\[
  \|V_k\|_p\ge c_p W_k^{1/p}.
\]
The same estimate holds for \(U_k\), after ignoring at most the two endpoint
row atoms.  On the other hand, (1) and the fact that \(u_k+v_k\) is supported
on edges with \(m\in J_k\) give a constant \(C_p\) such that
\[
  \|U_k+V_k\|_p\le C_p\frac{W_k^{1/p}}{L_k}.
\]
Thus
\[
  \frac{\|U_k+V_k\|_p}{\|U_k\|_p+\|V_k\|_p}\longrightarrow0
  \qquad(1\le p<\infty). \tag{3}
\]
For \(p=\infty\), the functions \(\theta_k\) have oscillation \(1\), whereas
\(\int u_k\,d\mu\to0\) and \(\int v_k\,d\mu\to0\).  Hence
\[
  \|U_k\|_\infty+\|V_k\|_\infty\ge 1
\]
for all large \(k\), while (2) tends to zero.  Therefore (3) also holds for
\(p=\infty\), with the usual interpretation.

Now fix \(p\in[1,\infty]\) and set
\[
  M=L^p_0(\A),\qquad N=L^p_0(\B).
\]
The addition map
\[
  T:M\oplus_1 N\to L^p(\Omega),\qquad T(f,g)=f+g,
\]
is bounded and injective by the lemma.  If \(M+N\) were closed, then \(T\)
would be a Banach-space isomorphism from \(M\oplus_1N\) onto \(M+N\), and its
inverse would be bounded.  Equivalently, there would exist \(\delta>0\) with
\[
  \|f+g\|_p\ge \delta(\|f\|_p+\|g\|_p)
  \qquad(f\in M,\ g\in N).
\]
The functions \(U_k,V_k\) violate this estimate by (3).  Therefore \(M+N\) is
not closed.

Since \(E=\left(\begin{smallmatrix}0&1\\1&0\end{smallmatrix}\right)\) has
spectral radius \(1\) and \(\psip(\A,\B)=0=1-1\), this proves the \(n=2\)
case of the conjecture.

\section*{Verification and limitations}

The proof is self-contained after the source definitions.  It proves only the
two-subalgebra case.  The full \(n\)-subalgebra conjecture remains open here,
especially for matrices of spectral radius one with all off-diagonal entries
strictly below one, such as the \(3\times3\) matrix with all off-diagonal
entries \(1/2\).

Bounded novelty search was performed on 2026-06-29 using the source title and
phrases involving \(\psip\), marginal subspaces, and nonclosed sums.  No later
paper resolving this exact \(n=2\) sharpness case was found in the search
results.

\begin{thebibliography}{9}

\bibitem{Feshchenko}
I. Feshchenko,
\emph{When is the sum of complemented subspaces complemented?},
arXiv:1606.08048.

\end{thebibliography}

\end{document}
